\documentclass[addpoints,12pt]{exam}

\usepackage{amsmath, amsfonts, amssymb}

% --- Header and Footer Setup ---
\pagestyle{headandfoot}
\firstpageheader{Comp Math}{Second Homework}{September 16, 2026}
\runningheader{Comp Math}{First Homework, Page \thepage\ of \numpages}{}
\firstpagefooter{}{}{}
\runningfooter{}{}{Page \thepage}

\begin{document}

\begin{center}
    \textbf{\Large Homework \#2} \\
    \vspace{0.1in}
    \textbf{\large Logic}
\end{center}

\vspace{0.3in}

\begin{questions}

% Q1
\question Analyze the logical forms of the following statements.
\begin{parts}
    \part Either John and Bill are both telling the truth, or neither of them is.
    \vspace{.1 in}

    let: 

    $J =$ John

    $B =$ Bill

    The logical form can be written out as: $(J \land B) \vee (\lnot J \land \lnot B) $

    \part I’ll have either fish or chicken, but I won’t have both fish and mashed potatoes.
    \vspace{.1 in}

    let: 

    $F =$ Fish

    $C =$ Chicken

    $M =$ Mashed Potatoes 

    The logical form can be written out as: $(F \vee C) \land \lnot(F \land M)$

    \part 3 is a common divisor of 6, 9, and 15
    \vspace{.1 in}

    The logical form can be written out as: $6,9,15 \in \{x | x \% 3 = 0\}$

\end{parts}


% Q2
\question Which of the following expressions are well-formed formulas?

\begin{parts}
\part $\lnot (\lnot P \land \lnot\lnot R)$, is well formed. While the double $\lnot$ is somewhat irksome its still correct 

\part $\lnot(P, Q, \vee R)$ No, as the commas make it not proper "Syntax" 

\part $P \land \lnot P$ Well formed, as it is clear and not ambigious. Unless parenthesis are required.

\part $(P \land Q)(P \vee R)$ Not well formed.
\end{parts}

\question Let T stand for the statement “Taxes will go up” and D for “The
deficit will go up.” What English sentences are represented by the
following formulas?

\begin{parts}
    \part $T \vee D$: Taxes will go up or The deficit will go up
    \part $¬(T \land D) \land \lnot(\lnot T \land \lnot D)$: 
    
    Taxes and Decificit will go up and Taxes and Deficit will both not not go up. 
    \part $(T \land \lnot D) \vee (D \land \lnot T)$:  

    Taxes and will go up and Deficit won't or Deficit will go up and Taxes wont. 
\end{parts}

\question Identify the premises and conclusions of the following deductive
arguments and analyze their logical forms. Do you think the reasoning
is valid? (Although you will have only your intuition to guide you in
answering this last question, in the next section we will develop some
techniques for determining the validity of arguments.)
\begin{parts}
    \part The main course will be either beef or fish. The vegetable will be
    either peas or corn. We will not have both fish as a main course and
    corn as a vegetable. Therefore, we will not have both beef as a main
    course and peas as a vegetable.
    \vspace{.1 in}

    False, as Beef and Peas are unconstrainted and are able to be a pair without a contradiction on the Fish and Corn constraint

    \part Either John or Bill is telling the truth. Either Sam or Bill is lying.
          Therefore, either John is telling the truth or Sam is lying.

    \vspace{.1 in}
    This is true, if John is true then it Makes Bill false and forces Sam to true, and When John is false, a similar chian of resoning makes Sam False 


\end{parts}

\large 1.2 Exersices 
\small

\question Make truth tables for the following formulas: 
\begin{parts}
    \part  $ \lnot [P \land (Q \vee \lnot P)] $
    \[
    \begin{array}{c|c|c|c|c|c}
    P & Q & \neg P & Q \vee \neg P & P \land (Q \vee \neg P) &
    \neg[P \land (Q \vee \neg P)] \\ \hline
    T & T & F & T & T & F \\
    T & F & F & F & F & T \\
    F & T & T & T & F & T \\
    F & F & T & T & F & T
    \end{array}
    \]

    \part $(P \vee Q) \land (\lnot P \vee R)$ 
    \[
    \begin{array}{c|c|c|c|c|c}
    P & Q & R & P \vee Q & \neg P \vee R &
    (P \vee Q) \land (\neg P \vee R) \\ \hline
    T & T & T & T & T & T \\
    T & T & F & T & F & F \\
    T & F & T & T & T & T \\
    T & F & F & T & F & F \\
    F & T & T & T & T & T \\
    F & T & F & T & T & T \\
    F & F & T & F & T & F \\
    F & F & F & F & T & F
    \end{array}
    \]
    
\end{parts}

\question In this exercise we will use the symbol + to mean exclusive or. In other words, P + Q means “P or Q, but not both.”

\begin{parts}
    \part Make a truth table for $P + Q$.

    \[
    \begin{array}{c|c|c}
    P & Q & P+Q \\ \hline
    T & T & F \\
    T & F & T \\
    F & T & T \\
    F & F & F
    \end{array}
    \]

    \part  Find a formula using only the connectives $\land, \vee$, and ¬ that is
equivalent to P + Q. Justify your answer with a truth table.

    \[
    \begin{array}{c|c|c|c|c|c}
    P & Q & P+Q & \neg Q & P\land\neg Q &
    \neg P\land Q \\ \hline
    T & T & F & F & F & F \\
    T & F & T & T & T & F \\
    F & T & T & F & F & T \\
    F & F & F & T & F & F
    \end{array}
    \]

    \[
    \begin{array}{c|c}
    (P\land\neg Q)\lor(\neg P\land Q) & P+Q \\ \hline
    F & F \\
    T & T \\
    T & T \\
    F & F
    \end{array}
    \]
    
\end{parts}

\question Find a formula using only the connectives $\land$ and $\lnot$ that is equivalent
to P $\vee$ Q. Justify your answer with a truth table.
    \begin{parts}
        \part when can make it with demorgans law
        \[
        \begin{array}{c|c|c|c|c|c}
        P & Q & P\vee Q & \neg P & \neg Q &
        \neg(\neg P\land\neg Q) \\ \hline
        T & T & T & F & F & T \\
        T & F & T & F & T & T \\
        F & T & T & T & F & T \\
        F & F & F & T & T & F
        \end{array}
        \]

    \end{parts}

\question Some mathematicians write P $|$ Q to mean “P and Q are not both true.” (This connective is called nand, and is used in the study of circuits in computer science.) 

\begin{parts}
    \part 
    \[
    \begin{array}{c|c|c|c}
    P & Q & P\land Q & P\mid Q=\neg(P\land Q) \\ \hline
    T & T & T & F \\
    T & F & F & T \\
    F & T & F & T \\
    F & F & F & T
    \end{array}
    \]
\end{parts}

\question Use truth tables to determine which of the following formulas are equivalent to each other
\begin{parts}
\part  $(P \land Q) \land (\lnot P \vee \lnot Q)$
    \[
    \begin{array}{c|c|c|c|c|c|c}
    P & Q & \neg P & \neg Q & P \lor Q & \neg P \lor \neg Q & (P \lor Q) \land (\neg P \lor \neg Q) \\
    \hline
    T & T & F & F & T & F & F \\
    T & F & F & T & T & T & T \\
    F & T & T & F & T & T & T \\
    F & F & T & T & F & T & F
    \end{array}
    \]
\part $\lnot P \vee Q$
    \[
    \begin{array}{c|c|c|c}
    P & Q & \neg P & \neg P\vee Q \\
    \hline
    T & T & F & T \\
    T & F & F & F \\
    F & T & T & T \\
    F & F & T & T
    \end{array}
    \]
\part $(P \vee \lnot Q) \land (Q \vee \lnot P)$ 
    \[
    \begin{array}{c|c|c|c|c|c|c}
    P & Q & \neg P & \neg Q & P\vee\neg Q & Q\vee\neg P & (P\vee\neg Q)\land(Q\vee\neg P) \\
    \hline
    T & T & F & F & T & T & T \\
    T & F & F & T & T & F & F \\
    F & T & T & F & F & T & F \\
    F & F & T & T & T & T & T
    \end{array}
    \]
\part $\lnot (P \vee Q)$
    \[
    \begin{array}{c|c|c|c}
    P & Q & P\vee Q & \neg(P\vee Q) \\
    \hline
    T & T & T & F \\
    T & F & T & F \\
    F & T & T & F \\
    F & F & F & T
    \end{array}
    \]
\part $(Q \land P) \vee \lnot P$
    \[
    \begin{array}{c|c|c|c|c}
    P & Q & Q\land P & \neg P & (Q\land P)\lor\neg P \\
    \hline
    T & T & T & F & T \\
    T & F & F & F & F \\
    F & T & F & T & T \\
    F & F & F & T & T
    \end{array}
    \]

(A) and (C) are equal and (B) and (E) and (D) is unique 
\end{parts}
\end{questions}
\end{document}